%! “This manuscript clarifies, within a single covariant vector--tensor framework, which core features of relativistic physics follow from Lorentzian geometry, which from variational structure, and which from minimal matter coupling.”
%!
%
%
\documentclass[a4paper,10pt]{letter}

\usepackage[T1]{fontenc}
\usepackage[utf8]{inputenc}
\usepackage{lmodern}
\usepackage[hidelinks]{hyperref}
\usepackage{microtype}
\usepackage[margin=1in]{geometry}

\signature{Omar Iskandarani\\
Independent Researcher, Groningen\\
The Netherlands\\
ORCID: 0009-0006-1686-3961\\
Email: \href{mailto:info@omariskandarani.com}{info@omariskandarani.com}}

\address{Omar Iskandarani\\
Vinkenstraat 86A\\
9713 TK Groningen\\
The Netherlands}

\date{\today}

\begin{document}
    
    \begin{letter}{Editors\\
    \textit{General Relativity and Gravitation}}

        \opening{Dear Editors,}
        
        Please consider the enclosed manuscript entitled
        \textbf{“Covariant Emergence of Special-Relativistic Kinematics and Einsteinian Gravitational Dynamics from a Minimal Vector--Tensor Action”}
        for publication in \textit{General Relativity and Gravitation}.
        
        This manuscript studies a focused foundational question within a minimal covariant vector--tensor framework equivalent to Einstein--\AE{}ther theory: which of the standard pillars of relativistic physics are structurally enforced by the action itself, and which depend on separate assumptions. In particular, the paper isolates three results within one common framework: local special-relativistic kinematics from Lorentzian metric structure, Einstein-form gravitational field equations from the Einstein--Hilbert variational principle, and geodesic free fall for minimally coupled structureless test bodies in the test-body regime.
        
        The contribution is not presented as new fundamental physics beyond the established Einstein--\AE{}ther literature. Rather, the paper offers a self-contained and logically explicit synthesis of known ingredients, showing that these three familiar results rest on distinct structural inputs and should not be conflated. The aim is therefore conceptual and foundational: to clarify the logical architecture of relativistic kinematics, gravitational dynamics, and weak-equivalence behavior within a single covariant action-based setting.
        
        The manuscript also includes a concise weak-field discussion that situates the framework within the standard post-Newtonian literature. In particular, it summarizes the inherited Einstein--\AE{}ther results for the preferred-frame sector and the role of the gravitational-wave speed constraint in restricting the viable parameter space.
        
        We believe this contribution is well suited to \textit{General Relativity and Gravitation}, given the journal’s longstanding interest in the conceptual structure of gravitational theory, generally covariant field frameworks, and careful foundational analysis of relativistic models.
        
        This manuscript is original, has not been published previously, and is not under consideration for publication elsewhere.
        
        Thank you for your time and consideration.
        
        \closing{Sincerely,}
    
    \end{letter}

\end{document}